\section{Strict cutpoints and numerical simulation}\label{sec:cutpoint65}
State cost for recognition and numerical precision are different
invariants. Recent results make it necessary to state the distinction
at theorem level, rather than relying on the words ``classical
simulation.''

Chen--Wu \cite[Definition~2.1 and Theorems~3.2, 4.4]{CWApril65}
define simulation by equality of strict-cutpoint languages, allowing
the simulator's cutpoint to change. Their August 2026 version gives
the exact worst-case cost $n^2+1$ for $n$-state one-way general
quantum automata, $n\ge2$. In \cite[Theorems~11, 14 and Corollary~15]{CWMay65}
they give order $cq^2$ for $c,q\ge2$ in a model with $c$ classical states and
a $q$-dimensional quantum register, and order $n^2$ for the
measure-once model. These are relevant positive results. The first
paper's current title is \emph{The State Cost of Classical Simulation
of One-Way General Quantum Finite Automata}; its version and exact
bound, rather than the earlier title or only its quadratic order,
are used here.

Their stochasticization preserves a sign, not the numerical
acceptance function. In particular,
\cite[Proposition~3.1]{CWApril65} contains an attenuation factor
$\delta\eta^{|w|}>0$ multiplying a zero-cutpoint linear output.
It can tend to zero exponentially without changing the recognized
language. The following elementary construction isolates this point.

\begin{proposition}[Attenuation preserves a strict cutpoint]
\label{prop:attenuation65}
Let a probabilistic automaton have acceptance function $f(w)$, and
fix $0<\beta<1$ and $0\le\lambda<1$. Adding two absorbing states
gives an automaton with acceptance function
\begin{equation}\label{eq:attenuation65}
             f_\beta(w)=\lambda+\beta^{|w|}(f(w)-\lambda).
\end{equation}
It recognizes the same strict-cutpoint language at $\lambda$.
If some length-$N$ word has $|f(w)-\lambda|\ge c>0$, its numerical
error on that word is at least $(1-\beta^N)c$.
\end{proposition}
\begin{proof}
On every letter, each old state follows its original stochastic row
with total probability $\beta$. The remaining mass goes to an
absorbing accepting state with probability $(1-\beta)\lambda$ and
to an absorbing rejecting state with probability
$(1-\beta)(1-\lambda)$. The old terminal functional is unchanged.
After $t$ letters the surviving old distribution is $\beta^t$
times the original distribution, while the absorbed accepting mass
is $\lambda(1-\beta^t)$. This proves \eqref{eq:attenuation65}.
The factor $\beta^{|w|}$ is positive, proving language equality.
Subtracting $f(w)$ proves the last assertion.
\end{proof}

A legal density-matrix approximation gives a numerical guarantee for
every supplied effect: if $0\le E\le I$ and $\rho,\sigma\in\mathcal D_d$,
then
\[
                 |\tr E(\rho-\sigma)|\le\tfrac12\|\rho-\sigma\|_1.
\]
Indeed the positive and negative parts of $\rho-\sigma$ have equal
trace, and an effect takes a value between zero and that trace on
either part. Therefore an approximation preserves a threshold
decision whenever its numerical error is strictly smaller than the
distance from the cutpoint. Without a margin, numerical approximation
alone does not give exact language equivalence.

The prepare--test and sign-rank methods in \cite{CWMay65} concern
finite prefix--suffix sign patterns. Such patterns are unchanged by
positive attenuation. Our entropy converse instead requires a
calibrated residual mean and a legal convex decoder; attenuation can
destroy that calibration. Conversely, its width lower bound cannot
be transferred to unrestricted cutpoint-language simulation. Constant
real linear dimension, fixed PFA state count and finite writable
bit-space are not identified with one another. No quantum language
separation, resolution of the Chen--Wu state-cost problem, or
unconditional lower bound for arbitrary noisy channels is claimed.
The new positive process construction is an upper theorem for a
different, explicitly specified numerical and transcript interface.
